%! Systems Involving Quadratics — Unit 2, Lesson 2.6 — Group Activity (Answer Key)
\documentclass[10pt]{article}
\usepackage{algebra2-article}
\usepackage{algebra2-key}   % pulls in -boxes; do NOT also load -boxes
\newcommand{\TallMath}[1]{$\displaystyle #1\rule[-1.4em]{0pt}{3.2em}$}
\newcommand{\lgrid}[4]{%
  \draw[step=1, linegray!40, very thin] (#1,#3) grid (#2,#4);
  \draw[->, semithick, charcoal] (#1,0)--(#2,0) node[below right, font=\scriptsize]{$x$};
  \draw[->, semithick, charcoal] (0,#3)--(0,#4) node[left, font=\scriptsize]{$y$};}

\begin{document}

\pageheader{Unit 2, Lesson 2.6}{Group Activity --- Answer Key}

\vspace{0.08in}

\begin{headlinebox}{forestbg}
\small\textbf{Where Do They Meet?} With your partner, remember the one big idea: a \textbf{solution} of a system is
a \textbf{point both graphs share}. To find it, set the two $y$-expressions \emph{equal} --- the system collapses
to a single quadratic you already know how to solve. Always finish by writing an \textbf{ordered pair} and checking
it in \emph{both} equations.
\end{headlinebox}

\vspace{0.06in}

\begin{tcolorbox}[colback=white, colframe=black!40, title=\textbf{Tier R --- Remediate},
    fonttitle=\bfseries, arc=1.5mm, left=3mm, right=3mm, top=2mm, bottom=2mm, breakable]
\textbf{1. Read the solutions off the graph.} The parabola $y=x^2-4$ meets the line $y=x-2$. Name the two
intersection points.
\begin{center}
\begin{tikzpicture}[scale=0.44]
  \lgrid{-3}{4}{-5}{4}
  \begin{scope}\clip (-3,-5) rectangle (4,4);
    \draw[very thick, forest, domain=-2.6:3.1, samples=60, smooth] plot (\x,{\x*\x-4});
    \draw[very thick, navy, domain=-3:4] plot (\x,{\x-2});
  \end{scope}
  \fill[charcoal] (2,0) circle (3pt); \fill[charcoal] (-1,-3) circle (3pt);
  \node[font=\scriptsize, forest] at (2.7,3.2){$y=x^2\!-\!4$};
  \node[font=\scriptsize, navy]   at (3.4,0.6){$y=x-2$};
\end{tikzpicture}
\end{center}
Solutions: \ans{$(2,0)$} and \ans{$(-1,-3)$}.

\textbf{2. How many solutions?} Count the intersection points in each picture (write $0$, $1$, or $2$).
\begin{center}
\begin{tikzpicture}[scale=0.34]
  \lgrid{-3}{3}{-4}{3}
  \begin{scope}\clip (-3,-4) rectangle (3,3);
    \draw[very thick, forest, domain=-2.6:2.6, samples=50, smooth] plot (\x,{\x*\x-3});
    \draw[very thick, navy] (-3,1)--(3,1);
  \end{scope}
  \fill[charcoal] (-2,1) circle (2.4pt); \fill[charcoal] (2,1) circle (2.4pt);
  \node[font=\scriptsize\bfseries, charcoal] at (0,-5.0){A: \ans{$2$}};
\end{tikzpicture}
\hspace{0.5cm}
\begin{tikzpicture}[scale=0.34]
  \lgrid{-2}{2}{0}{4}
  \begin{scope}\clip (-2,0) rectangle (2,4);
    \draw[very thick, forest, domain=-1.7:1.7, samples=50, smooth] plot (\x,{\x*\x+1});
    \draw[very thick, navy] (-2,1)--(2,1);
  \end{scope}
  \fill[charcoal] (0,1) circle (2.4pt);
  \node[font=\scriptsize\bfseries, charcoal] at (0,-1.0){B: \ans{$1$}};
\end{tikzpicture}
\hspace{0.5cm}
\begin{tikzpicture}[scale=0.34]
  \lgrid{-2}{2}{-2}{4}
  \begin{scope}\clip (-2,-2) rectangle (2,4);
    \draw[very thick, forest, domain=-1.7:1.7, samples=50, smooth] plot (\x,{\x*\x+1});
    \draw[very thick, navy] (-2,-1)--(2,-1);
  \end{scope}
  \node[font=\scriptsize\bfseries, charcoal] at (0,-3.0){C: \ans{$0$}};
\end{tikzpicture}
\end{center}

\textbf{3. Is it a solution?} A pair is a solution only if it works in \emph{both} equations. Check each against
$\ y=x^2-3\ $ and $\ y=x-3$ (write yes/no, and say why).
\begin{itemize}\setlength{\itemsep}{4pt}
  \item $(1,-2)$: \ parabola? \ans{yes} line? \ans{yes} $\Rightarrow$ solution? \ans{yes}
  \item $(2,1)$: \ parabola? \ans{yes} line? \ans{no} $\Rightarrow$ solution? \ans{no (on the parabola only)}
\end{itemize}
\end{tcolorbox}

\begin{tcolorbox}[colback=white, colframe=black!40, title=\textbf{Tier A --- Approaching Proficiency},
    fonttitle=\bfseries, arc=1.5mm, left=3mm, right=3mm, top=2mm, bottom=2mm, breakable]
\textbf{1. Solve by substitution} (linear--quadratic). Set the $y$'s equal, solve, back-substitute, write pairs.
\[ y=x^2+2x-1,\ \ y=x+1:\ \ \ans{$(-2,-1)$ and $(1,2)$} \]
\[ y=x^2-2x+2,\ \ y=-x+4:\ \ \ans{$(2,2)$ and $(-1,5)$} \]

\textbf{2. Quadratic--quadratic.} Set the expressions equal.
\[ y=x^2+1,\ \ y=-x^2+9:\ \ \ans{$(-2,5)$ and $(2,5)$} \]

\textbf{3. It won't factor --- use the quadratic formula.} Give the $x$-coordinates in simplest radical form.
\[ y=x^2-2,\ \ y=2x-1:\ \ x=\ans{$1\pm\sqrt{2}$} \]
\end{tcolorbox}

\begin{tcolorbox}[colback=white, colframe=black!40, title=\textbf{Tier E --- Extension},
    fonttitle=\bfseries, arc=1.5mm, left=3mm, right=3mm, top=2mm, bottom=2mm, breakable]
\textbf{1. Model (break-even).} A shop's revenue is $y=-x^2+14x$ and its cost is $y=4x+16$ (hundreds of dollars,
$x$ in hundreds of units). Find both break-even points, then say where the shop makes a profit.
\[ x=\ans{$2,\ 8$} \qquad\text{points: } \ans{$(2,24),(8,48)$} \qquad\text{profit when } \ans{$2<x<8$} \]

\textbf{2. Make it tangent (discriminant as a dial).} For what value of $c$ is the line $y=2x+c$ \emph{tangent} to
the parabola $y=x^2$ (exactly one solution)? \emph{Hint: set equal, then force $b^2-4ac=0$.}
\[ c=\ans{$-1$} \]

\textbf{3. Verify.} Show $(2,5)$ is a solution of the Tier A \#2 system by checking it in \emph{both} equations.
\ansline{$y=x^2+1:\ 2^2+1=5$ \checkmark; \ $y=-x^2+9:\ -(2^2)+9=5$ \checkmark. Both hold, so $(2,5)$ is a solution.}
\end{tcolorbox}

\end{document}
